\subsection{\label{sec:hmcdiag}\index{HMC diagnostics}HMC diagnostics}

The HMC diagnostic suite contains two tests: molecular-dynamics
reversibility and a finite-difference check of the production force.  The
implementation is in
\texttt{Code/CLGTest/Tests/TestHmcDiagnostics.cpp}; its test configurations
are in \texttt{Bin/Debug/TestSuit\_HmcDiagnostics.yaml}.

\subsubsection{\label{sec:hmcdiag_state}Gauge-dependent state}

Let $U$ denote the thin gauge field supplied to an action.  Before evaluating
the force or energy at a particular value of $U$, the diagnostic rebuilds all
objects derived from that same $U$ in the production order:
\begin{enumerate}
  \item update the registered cache before gauge updating;
  \item update the registered cache before smearing;
  \item construct the registered smeared or effective gauge field, if the
        smearing object is enabled during an update;
  \item update the registered cache after smearing.
\end{enumerate}
Thus an energy evaluated at a perturbed gauge field never uses a smeared field,
staple, plaquette, or $F_{\mu\nu}$ cache made from a different gauge field.

For a fermion action, the diagnostic first performs this refresh at the
unperturbed field $U_0$ and then calls \texttt{PrepareForHMC} once.  Denote the
resulting pseudofermion field by $\phi_0$.  The same $\phi_0$ is used for the
production-force calculation and for every positive and negative perturbation
defined below.  In particular, $\phi_0$ is not regenerated when the link,
direction, or finite-difference step size changes.

\subsubsection{\label{sec:hmcdiag_force}Production-force convention}

For each link $l$, \texttt{CalculateForce} adds a raw matrix $Y_l$ to the raw
force field.  The force actually used by the integrator is
\begin{equation}
F_l={\rm TA}\!\left(U_{0,l}Y_l^\dagger\right),
\label{eq.hmcdiag.force}
\end{equation}
where $U_{0,l}$ is link $l$ of the unperturbed gauge field and
\begin{equation}
{\rm TA}(M)
=\frac{M-M^\dagger}{2}
-\frac{1}{3}{\rm tr}\!\left(\frac{M-M^\dagger}{2}\right)\mathbb I.
\label{eq.hmcdiag.ta}
\end{equation}
Here $\mathbb I$ is the $3\times3$ identity matrix and ${\rm tr}$ is the
color-matrix trace.
In the code, Eq.~(\ref{eq.hmcdiag.force}) is implemented by
\texttt{LeftMul(U0,FALSE,TRUE)} followed by \texttt{TA()}.  Directions fixed
by \texttt{BindDir} are then set to zero in $F$.

For two gauge-shaped matrix fields $A$ and $B$, define the real Frobenius inner
product
\begin{equation}
\langle A,B\rangle
=\mathop{\rm Re}\sum_l {\rm tr}\!\left(A_l^\dagger B_l\right).
\label{eq.hmcdiag.inner}
\end{equation}
This is the real part returned by \texttt{A.Dot(B)}.  The CLGLib force
normalization predicts the action derivative
\begin{equation}
D_{\rm force}=-2\langle F,X\rangle.
\label{eq.hmcdiag.fdforcefield}
\end{equation}
The perturbation field $X$ used below is nonzero on only one link $l$.
Consequently Eq.~(\ref{eq.hmcdiag.fdforcefield}) reduces to
\begin{equation}
D_{\rm force}
=-2\mathop{\rm Re}{\rm tr}\!\left(F_l^\dagger X_l\right),
\label{eq.hmcdiag.fdforce}
\end{equation}
which is calculated by
\texttt{-2.0*pForce->Dot(pX).x}.

\subsubsection{\label{sec:hmcdiag_fd}Finite-difference force test}

The diagnostic tests each continuous action separately.  Fix one such action,
its unperturbed thin gauge field $U_0$, and, for a fermion action, the
pseudofermion field $\phi_0$ defined above.  The action value written below is
therefore
\begin{equation}
S[U]\equiv S[U;\phi_0]
\end{equation}
for a fermion action; for a gauge action there is no pseudofermion argument.

Choose a link $l$ which is neither Dirichlet nor fixed by \texttt{BindDir}.
Choose a $3\times3$ matrix $X_l$ satisfying
\begin{equation}
X_l^\dagger=-X_l,\qquad
{\rm tr}(X_l)=0,\qquad
\lVert X_l\rVert_{\rm F}=1,
\label{eq.hmcdiag.x}
\end{equation}
where
\begin{equation}
\lVert X_l\rVert_{\rm F}^2
={\rm tr}(X_l^\dagger X_l).
\end{equation}
Define the gauge-shaped field $X$ by setting its value to $X_l$ on link $l$
and to zero on every other link.  This definition makes
$\langle F,X\rangle$ in Eq.~(\ref{eq.hmcdiag.fdforcefield}) depend only on the
selected link.

Let $\Delta>0$ be the YAML parameter \texttt{Delta}, and let $N_h$ be the
integer YAML parameter \texttt{DeltaLevels}.  For
$k=0,1,\ldots,N_h-1$, define
\begin{equation}
h_k=\frac{\Delta}{2^k}.
\label{eq.hmcdiag.hk}
\end{equation}
For each $k$, make two independent copies of $U_0$.  Change only link $l$ and
define
\begin{align}
U_{+,k}(l)&=\exp(h_kX_l)U_0(l),
&U_{+,k}(l')&=U_0(l')\quad(l'\ne l),
\label{eq.hmcdiag.uplus}\\
U_{-,k}(l)&=\exp(-h_kX_l)U_0(l),
&U_{-,k}(l')&=U_0(l')\quad(l'\ne l).
\label{eq.hmcdiag.uminus}
\end{align}
The multiplication is on the left.  In the implementation it is performed by
\texttt{pX->ExpMult(+h,pWorkGauge)} and
\texttt{pX->ExpMult(-h,pWorkGauge)}.  Because $X_l$ is traceless and
anti-Hermitian, both exponentials are in $SU(3)$, so the perturbed link remains
in $SU(3)$.

After rebuilding all gauge-dependent data separately from $U_{+,k}$ and
$U_{-,k}$, calculate
\begin{equation}
S_{+,k}=S[U_{+,k}],\qquad S_{-,k}=S[U_{-,k}].
\label{eq.hmcdiag.spm}
\end{equation}
The central finite-difference estimate of the directional derivative at
$U_0$ is
\begin{equation}
D_{{\rm FD},k}
=\frac{S_{+,k}-S_{-,k}}{2h_k}.
\label{eq.hmcdiag.fd}
\end{equation}
The production-force prediction $D_{\rm force}$ in
Eq.~(\ref{eq.hmcdiag.fdforce}) is calculated once at $U_0$ and is the same for
all $k$.

For every level define the absolute disagreement
\begin{equation}
e_k=\left|D_{{\rm FD},k}-D_{\rm force}\right|,
\label{eq.hmcdiag.error}
\end{equation}
the comparison scale
\begin{equation}
q_k=\max\!\left(\left|D_{{\rm FD},k}\right|,
                 \left|D_{\rm force}\right|\right),
\label{eq.hmcdiag.scale}
\end{equation}
and the allowed error
\begin{equation}
b_k=\epsilon_{\rm abs}+\epsilon_{\rm rel}q_k,
\label{eq.hmcdiag.bound}
\end{equation}
where $\epsilon_{\rm abs}$ and $\epsilon_{\rm rel}$ are the YAML parameters
\texttt{TolAbs} and \texttt{TolRel}, respectively.

A link-and-direction trial fails if any required value
$S_{+,k}$, $S_{-,k}$, $D_{{\rm FD},k}$, $D_{\rm force}$, $e_k$, or $b_k$ is
not finite.  Otherwise the trial passes exactly when
\begin{equation}
\text{there exists }k\in\{0,\ldots,N_h-1\}\text{ such that }e_k\leq b_k.
\label{eq.hmcdiag.pass}
\end{equation}
There is no additional plateau, monotonic-decrease, convergence-rate, or
adjacent-level requirement.  The code reports all levels and identifies the
level with the smallest $e_k$, but the smallest-error level has no separate
acceptance rule beyond Eq.~(\ref{eq.hmcdiag.pass}).  Every sampled direction
on every sampled link must pass, and every selected continuous action is
tested independently.

The central difference in Eq.~(\ref{eq.hmcdiag.fd}) has an $O(h_k^2)$
truncation error in exact arithmetic.  Several values of $h_k$ are used because
decreasing $h_k$ reduces this truncation error but can amplify subtraction,
round-off, and fermion-solver errors.  This observation explains the step-size
sweep; it is not an extra pass condition.

\subsubsection{\label{sec:hmcdiag_rev}Reversibility}

Reversibility is tested using one prepared Hamiltonian state.  The test calls
\texttt{Prepare(FALSE,0)} exactly once, which prepares the action state and
generates the initial momentum.  Let the saved integrator working gauge and
momentum immediately after that call be $U_0$ and $P_0$.  The test performs
\begin{equation}
(U_0,P_0)\longrightarrow(U_1,P_1),
\end{equation}
negates the momentum, and then performs a second evolution without another
call to \texttt{Prepare}:
\begin{equation}
(U_1,-P_1)\longrightarrow(U_2,P_2).
\end{equation}
The same integrator, step size, step count, action state, and boundary
conditions are therefore used in both evolutions.  Reversibility requires
\begin{equation}
U_2\simeq U_0,\qquad P_2\simeq -P_0.
\end{equation}
The diagnostic reports average and maximum Frobenius norms of $U_2-U_0$ and
$P_2+P_0$ and compares them with the configured tolerances.

The two calls to \texttt{Evaluate()} modify the integrator working fields but
do not copy them to the lattice gauge field.  Reversibility checks the complete
molecular-dynamics map, but it does not prove that the force is correct: an
incorrect force can still be used reversibly in both evolutions.
